Arran de l'anunci fet el desembre del 2022 pels Estats Units que s'havia aconseguit per primera vegada la fusió nuclear amb un guany net d'energia, alguns mitjans de comunicació han publicat que un got d'aigua pot produir l'energia que consumirà una família de quatre membres durant tota la vida. Per una altra banda, a dins del Sol hi ha una pressió i una temperatura tan elevades que els àtoms d'hidrogen es fusionen i es transformen en heli. El principal procés de fusió nuclear que té lloc al Sol és la cadena protó-protó. El balanç global d'aquest procés és que quatre protons s'uneixen per formar un nucli d'heli.

\begin{apartat}{1,25}
Calculeu l'energia que s'allibera en la cadena protó-protó quan es forma un nucli d'heli. Tenint en compte que un got d'aigua conté, aproximadament, 17~mol d'aigua, quanta energia es pot extreure de l'hidrogen que hi ha a l'aigua d'un got mitjançant la cadena protó-protó?
\end{apartat}

\begin{apartat}{1,25}
Malauradament, a la Terra no es poden assolir les condicions de temperatura i pressió que hi ha al Sol, per això, en els reactors de fusió es fan servir els isòtops de l'hidrogen: el deuteri (${}^{2}_{1}\mathrm{H}$) i el triti (${}^{3}_{1}\mathrm{H}$). L'abundància relativa del deuteri és d'un 0,001~\%, mentre que la del triti és pràcticament nul·la (a la Terra hi ha uns 20~kg de triti natural).\newline
\hspace*{1.5em}Per a generar triti, s'utilitza liti 6 (${}^{6}_{3}\mathrm{Li}$) obtingut a partir de reactors nuclears de fissió. El triti s'obté a còpia de bombardejar nuclis de liti 6 amb neutrons. Escriviu la reacció nuclear sabent que el resultat és la formació de triti i partícules alfa. S'ha afirmat que la fusió nuclear és una energia neta perquè la reacció de la cadena protó-protó no genera residus radioactius. Ara bé, en el procés de fusió del deuteri i del triti s'alliberen neutrons amb una energia capaç de fer tornar radioactius els materials circumdants. Digueu si la fusió nuclear és una font d'energia neta i inesgotable i justifiqueu la resposta.
\end{apartat}

\dades{%
  $c = 3,00\times 10^{8}\ \mathrm{m\,s^{-1}}$.\\
  $N_\mathrm{A} = 6,022\times 10^{23}$.\\
  Masses nuclears (en kg):\\[3pt]
  \begin{tabular}{|c|c|}
    \hline
    \textit{Protó} & \textit{Nucli d'heli}\\
    \hline
    $1,672\,621\,92\times 10^{-27}$ & $6,642\,835\,33\times 10^{-27}$\\
    \hline
  \end{tabular}\\[3pt]
  El consum anual mitjà d'electricitat d'una persona a Catalunya és d'$1,22\times 10^{10}\ \mathrm{J}$.}
